\paragraph{Two Groups---Threshold Policies}

We now discuss the effects of a threshold policy 
when candidates belong to two groups, $G_1$ and $G_2$ 
whose skill level is distributed identically,
but whose tests are characterized by different noise levels.
Without loss of generality, we assume that $\eta_1<\eta_2$, 
where $\eta_i$ is the probability 
that a test result of a candidate from $G_i$ is different from his skill level.
To begin, we note the fundamental irreconcilability of equalizing either the false positive or the false negative rates across groups with subjecting candidates to the same policy.

\begin{theorem}[Impossibility result]\label{thm:stochasticDominance}
When noise levels differ between two groups with identical skill level distribution, 
a single Threshold Policy $\pi$ 
(with the same number of tests $\tau$ 
and the same threshold $\theta$ for both groups) 
cannot have equality in either the false negative rates 
or in the false positive rates across the groups. 
Particularly, there is a higher false positive rate in the noisier group, 
as an unskilled candidate from $G_2$ is more likely to be accepted 
by the threshold policy than an unskilled candidate from $G_1$:
\begin{equation*}
\textnormal{FPR}_{\theta,\tau}^{\eta_1}=
\Pr_{\eta_1}[\pi(\hat{y}_{i,1},\dots,\hat{y}_{i,\tau})=1|y_i=0]<\Pr_{\eta_2}[\pi(\hat{y}_{i,1},\dots,\hat{y}_{i,\tau})=1|y_i=0]
=\textnormal{FPR}_{\theta,\tau}^{\eta_2},
\end{equation*}
and also a higher false negative rate, 
as a skilled candidate from $G_2$ is more likely to be rejected 
than a skilled candidate from $G_1$:
\begin{equation*}
\textnormal{FNR}_{\theta,\tau}^{\eta_1}=
\Pr_{\eta_1}[\pi(\hat{y}_{i,1},\dots,\hat{y}_{i,\tau})=0|y_i=+1]<\Pr_{\eta_2}[\pi(\hat{y}_{i,1},\dots,\hat{y}_{i,\tau})=0|y_i=+1]=\text{FNR}_{\theta,\tau}^{\eta_2}.
\end{equation*}
\end{theorem}
\textbf{Connection to Economics Literature }%
 Aigner and Cain \cite{aigner1977statistical} discuss 
a similar case under a Gaussian screening model
where the variance (noise level) of the single test
differs across the two groups. 
Similarly, they note that qualified candidates fare worse in the noisy group
but that unqualified candidates fare better in the noisier group.
Our work differs from theirs in that 
we consider the effect of multiple tests
and the ability to optimize over the number of tests.

\paragraph{Two Groups--Dynamic policy}
We now consider the (dynamic) hiring policy 
in the setting when employees belong to two groups, 
$G_1$ and $G_2$ with identically-distributed skills 
but different noise levels $\eta_1<\eta_2$.
We note that there are two ingredients 
that explain the differences among the groups:
(i) The step size, $\log \left(\frac{\Pr[\hat{y}_{i,j}=+1|y_i=+1]}{\Pr[\hat{y}_{i,j}=+1|y_i=0]}\right)=\log \left(\frac{1-\eta}{\eta} \right)$ of $G_2$ (the noisier group) 
is smaller than the step size of $G_1$. 
Thus these candidates must typically pass more tests 
before they are accepted;
and (ii)
Skilled candidates in group $G_2$ 
exhibit less drift to the right
(they have a higher probability of failing a test).
Consequently, when an employer (rationally)
sets $\epsilon' = p$ for all groups,
a skilled candidate from $G_2$ 
is more likely to be fail a test in step $1$,
at which point the dynamic policy summarily rejects them.
These two facts explain both the higher false negative rates for $G_2$ 
and the longer expected duration until acceptance.
By setting $\epsilon' < p$ for members of the noisier group, we can equalize false negative rates.
Precisely, setting $\epsilon' = \frac{\eta_1}{\eta_2}p$
achieves the desired parity.
The cost of this intervention is that it requires more tests for candidates from the noisier group.
Here, our random walk analysis can be leveraged to determine exactly how many more.
Once again, we cannot provide equality across the groups in all desired ways---the 
same acceptance criterion, the same expected number of tests,
and the same false negative rates between groups---with the noise differs across groups.